%---------------------------------------------------------------------
\begin{frame}{}

    \begin{center}
        \Large \textcolor{maroon}{The Grossman model (1972)}
        \end{center}
\end{frame}
%---------------------------------------------------------------------
\begin{frame}{}

    Picture first page of article
\end{frame}
%---------------------------------------------------------------------
\begin{frame}{What is \textit{health}?}

    \begin{itemize}
        \item Something we enjoy: a \newbold{consumption good}
        \pause
        \item Something we invest in: a \newbold{capital good}
        \pause
        \item Something that affects our ability to work: an \newbold{input} in production of productive time

    \end{itemize}
\end{frame}
%---------------------------------------------------------------------
\begin{frame}{Single-period utility}

    Utility
    \begin{equation*}
        U = U\left( H_t, Z_t \right)
    \end{equation*}
    where
    \begin{itemize}
        \item $H_t$ is the level of health at time $t$
        \item $Z_t$ is a composite consumption good at time $t$
    \end{itemize}


\end{frame}
%---------------------------------------------------------------------
\begin{frame}{Single-period time constraint}

    An individual's activity:
    \begin{itemize}
        \item Spends time working (and earns a wage $w$ from it): $T^w$
        \item Spends time consuming the composite good $Z$: $T^Z$
        \item Spends time investing in health $H$: $T^H$
        \item If sick, needs to rest in bed and cannot do anything else: $T^s$
    \end{itemize}
    \vspace{3mm}

    Time constraint:
    \begin{equation*}
        \Theta = T^w + T^Z + T^H + T^s
    \end{equation*}
    where $\Theta$ is total available time in a period.

\end{frame}
%---------------------------------------------------------------------
\begin{frame}{Stock versus flow}

    \begin{equation*}
        U = U\left( H_t, Z_t \right)
    \end{equation*}

    Note:
    \begin{itemize}
        \item $Z_t$ is a \textit{flow}
        \item $H_t$ is a \newbold{stock}
    \end{itemize}
    $\implies$ How do you produce them?

\end{frame}
%---------------------------------------------------------------------
\begin{frame}{Production}

    Production functions:
    \begin{itemize}
        \item Health:
        \begin{equation*}
            H_t = H\left(H_{t-1}, T_t^H, M_t \right)
        \end{equation*}
        where
        \begin{itemize}
            \item $H_{t-1}$ is health stock from previous period
            \item $T_t^H$ is time invested in health this period
            \item $M_t$ are market inputs invested in health (e.g., vaccine, treadmill)
        \end{itemize}
        \pause
        \item Composite good:
        \begin{equation*}
            Z_t = Z\left( T_t^Z, J_t \right)
        \end{equation*}
        where
        \begin{itemize}
            \item $T_t^Z$ is time invested in consuming $Z$
            \item $J_t$ are market inputs invested in consuming $Z$ (e.g., food, video games)
        \end{itemize}
    \end{itemize}
    $\implies$ Health is a stock: previous period's stock enters the production function//
    $\implies$ You need time to get both consumption and health!
\end{frame}
%---------------------------------------------------------------------
\begin{frame}{Another constraint: the budget constraint}

    Time constraint \textit{(which we already saw)}:
    \begin{equation*}
        \Theta = T^w + T^Z + T^H + T^s
    \end{equation*}
    \vspace{3mm}
    \pause

    \only<1>{Budget constraint:
    \begin{equation*}
        p_M M_t + p_J J_t \leq w T^w
    \end{equation*}
    where
    \begin{itemize}
        \item $p_M$ is the price of market inputs for health
        \item $p_J$ is the price of market inputs for consumption
        \item $w$ is the wage rate
    \end{itemize}
    }

    \only<2>{Budget constraint:
    \begin{equation*}
        p_M M_t + p_J J_t \textcolor{red}{=} w T^w
    \end{equation*}
    (no \textcolor{red}{saving/borrowing})}

    \only<3>{Budget constraint:
    \begin{equation*}
        p_M M_t + p_J J_t = w T^w
    \end{equation*}
    \vspace{3mm}

    Note: the two constraints are \newbold{not independent}.
    }

\end{frame}
%---------------------------------------------------------------------
\begin{frame}{Health and time sick}

    Productive time:
    \begin{align*}
        T^p &= \Theta - T^s\\
            &= T^w + T^Z + T^H
    \end{align*}
    \vspace{3mm}

    Time sick is a function of health stock:
    ADD graphics
    \pause
    \begin{itemize}
        \item Better health $\implies$ less time sick
        \item \newbold{Diminishing marginal returns}
        \item Where is \textit{death}? $H_min$
    \end{itemize}


\end{frame}
%---------------------------------------------------------------------
\begin{frame}{Production possibility frontier of $Z$ and $H$}

    ADD FIGURE 3.2 and 3.3 from book


\end{frame}
%---------------------------------------------------------------------
\begin{frame}{Single-period optimal production}

    ADD FIGURE 3.4 + exotic indifference curves


\end{frame}
%---------------------------------------------------------------------
\begin{frame}{Optimal time working?}

    Fix $T^H$ for a minute

    ADD FIGURE 3.8 from book

    Now assume $H_t$ is larger:
    ADD figure 3.9


\end{frame}
%---------------------------------------------------------------------
\begin{frame}{Multi-period model}

    Assume multi-period utility as
    \begin{align*}
        U &= U(H_0, Z_0) + \delta U(H_1, Z_1) + \delta^2 U(H_2, Z_2) + \ldots + \delta^T U(H_T, Z_T)\\
          &=\sum_{t=0}^{T} \delta^t U\left( H_t, Z_t \right)
    \end{align*}
    where $\delta$ is the discount factor.
    \vspace{3mm}
    \pause

    Health as an investment good: introduce \textit{aging}
    \begin{equation*}
        H_t = H\left((1 - \gamma)H_{t-1}, T_t^H, M_t \right)
    \end{equation*}
    where $\gamma$ is the \newbold{depreciation rate} of health.

\end{frame}
%---------------------------------------------------------------------
\begin{frame}{Lifetime return to a marginal investment in health}

    ADD Figure 3.10

    ADD Figure 3.11

\end{frame}
%---------------------------------------------------------------------
\begin{frame}{Aging}

    How to think about aging? change $\gamma$

    ADD Figure 3.13

\end{frame}
%---------------------------------------------------------------------
\begin{frame}{Relationship between SES and health}

    Take some graphs from slides of Vicky

\end{frame}
%---------------------------------------------------------------------
\begin{frame}{Using the Grossman model to think about the SES gradient in health}

    ADD Figure 3.12

\end{frame}
%---------------------------------------------------------------------
\begin{frame}{}

    \begin{center}
        \Large \textcolor{maroon}{Let's go to worksheet 1}
        \end{center}
\end{frame}
%---------------------------------------------------------------------
%---------------------------------------------------------------------
\begin{frame}{}
    \begin{center}
    \Large \textcolor{maroon}{Wraping up}
    \end{center}

\end{frame}
%---------------------------------------------------------------------
\begin{frame}{What we covered today}
    
    \newbold{Demand for health: the Grossman model (1972)}
    \begin{itemize}
        \item Aging
        \item SES and health
    \end{itemize}
\end{frame}
%---------------------------------------------------------------------
\begin{frame}{To next class}
    \begin{itemize}
        \item \newbold{TO DO}
        \begin{enumerate}
            \item Finish worksheet 1 $\implies$ \textit{questions on it?}
            \item Review past material $\implies$ \textit{questions on it?}
        \end{enumerate}
        \vspace{5mm}

        \item Program for next class
        \begin{itemize}
            \item Health insurance
        \end{itemize}
    \end{itemize}
\end{frame}
%---------------------------------------------------------------------